% ============================================================
%  Chapter 7: Conclusion
% ============================================================
%
%  TARGET: ~300 words. Tight. Write LAST.
%
%  HARD RULES:
%   - Do not introduce new information.
%   - Do not repeat Chapter 5 numbers verbatim -- the
%     introduction has them; here just refer.
%   - Keep future work to 3-4 items max.
%
% ============================================================
%
%  TARGET STRUCTURE (no subsections needed at this length, but
%  keep them if it helps your template):
%
%  ¶1 Summary (~120 w)
%       One sentence on the problem (sparse-view CT is
%       ill-posed). One on the approach (faithful reproduction
%       of TFPnP in LION with DRUNet as the off-the-shelf
%       denoiser, trained on LIDC-IDRI). One on the headline
%       result (TFPnP beats Fixed PnP-ADMM by ~X dB on the 389-
%       image test set across three noise levels, with the
%       policy learning interpretable image-specific schedules).
%
%  ¶2 Contributions (~80 w, compact list)
%       (i) In-domain CT reproduction of TFPnP with
%           sigma/mu calibration for mu-valued inputs.
%       (ii) DRUNet integration validating the off-the-shelf
%            generalisation claim with a non-IRCNN denoiser.
%       (iii) Empirical finding: depth-per-image dominates
%            data volume in RL-tuned PnP training.
%       (iv) LION toolbox integration (engineering).
%
%  ¶3 Future work (~100 w, 3-4 bullets in prose)
%       - LPD and FBPconv trained baselines for full
%         like-for-like learned-method comparison.
%       - DRUNet fine-tuned on CT, or a CT-specific denoiser
%         (e.g. ACR/AR) -- testing whether the residual
%         over-smoothing closes the perceptual gap.
%       - Cone-beam geometry (LION's tomosipo backend supports
%         this natively, unlike the paper's TorchRadon).
%       - Realistic noise model (Poisson on photon counts) for
%         clinical applicability.
%
% ============================================================
%  FRAMING TO USE:
%   "faithful reproduction" not "reproduction of"
%   "in-domain advantage" when discussing paper comparisons
%   "depth-per-image" -- your novel finding, name it
%
%  FRAMING TO AVOID:
%   Overclaiming the LION integration as research contribution.
%   Claiming the work surpasses the paper (in-domain caveat).
%   Vague platitudes like "results demonstrate the power of RL".
% ============================================================

\chapter{Conclusion}
\label{chap:conclusion}

% TODO: write after Ch.5 and Ch.6 are finalised so headline numbers
% can be quoted accurately.